% --- CODIFICACIÓN Y FUENTES ---
\usepackage[utf8]{inputenc}
\usepackage[T1]{fontenc}
% Estilo Del Valle: texto en Times, pero matemáticas en Computer Modern.
% (mathptmx pondría también las matemáticas en Times; aquí solo cambiamos el texto.)
\renewcommand{\rmdefault}{ptm}

% --- IDIOMA ---
\usepackage[english,spanish,es-nodecimaldot]{babel}

% --- GEOMETRÍA DE PÁGINA ---
\usepackage[a4paper,
            top=3cm,
            bottom=3cm,
            left=3.5cm,
            right=2.5cm,
            headheight=15pt
           ]{geometry}

% --- ENCABEZADOS / PIE (estilo Del Valle) ---
% Sin encabezados ni reglas; número de página centrado en el pie.
\usepackage{fancyhdr}
\pagestyle{fancy}
\fancyhf{} % Limpia todo
\fancyfoot[C]{\thepage} % Número de página centrado abajo
\renewcommand{\headrulewidth}{0pt}
\renewcommand{\footrulewidth}{0pt}

% El mismo aspecto en las páginas de inicio de capítulo (estilo plain)
\fancypagestyle{plain}{%
  \fancyhf{}%
  \fancyfoot[C]{\thepage}%
  \renewcommand{\headrulewidth}{0pt}%
  \renewcommand{\footrulewidth}{0pt}%
}

\setlength{\headheight}{15.2pt}

% --- MATEMÁTICAS ---
\usepackage{amsmath}
\usepackage{amsfonts}
\usepackage{amssymb}

% --- UNIDADES ---
\usepackage{siunitx}
\sisetup{
    range-phrase = {--},
    range-units = single,
    separate-uncertainty=true,
    detect-all
}

% --- IMÁGENES Y FIGURAS ---
\usepackage{graphicx}
\graphicspath{{./images/}}
\usepackage{caption}
\usepackage[font=small,labelfont=bf,labelsep=colon]{caption}
\usepackage{subcaption}
\usepackage{float}      
\usepackage{eso-pic}    

% --- TABLAS ---
\usepackage{tabularx}
\usepackage{booktabs}
\usepackage{multirow}
\usepackage{array}

% --- LISTAS Y FORMATO ---
\usepackage{multicol}   % Lista de acrónimos a dos columnas (estilo Del Valle)
\usepackage{enumitem}
\usepackage{setspace}
% \usepackage{parskip}    
\usepackage{microtype}
\usepackage{titlesec}

% Configuración de Títulos de Capítulos (estilo Del Valle:
% "Capítulo N" en negrita y, debajo, el título en negrita más grande)
\titleformat{\chapter}[display]
  {\normalfont\huge\bfseries}{\chaptertitlename\ \thechapter}{14pt}{\Huge}
\titlespacing*{\chapter}{0pt}{50pt}{30pt}

% --- CITAS Y BIBLIOGRAFÍA ---
\usepackage[square,numbers,sort&compress]{natbib}

% --- HIPERVÍNCULOS (Al final) ---
\usepackage{xcolor}
% Paleta estilo Del Valle: enlaces internos en teal-azul (muestreado del índice
% original, RGB 0/96/136 ~ #006088), citas en rojo oscuro.
\definecolor{linkblue}{rgb}{0.000, 0.376, 0.533}
\definecolor{citered}{rgb}{0.55, 0.0, 0.0}

\usepackage{hyperref}
\hypersetup{
    colorlinks=true,
    linktoc=page,        % Estilo Del Valle: en el índice solo se colorea el número de página (los títulos quedan en negro)
    linkcolor=linkblue,
    citecolor=citered,
    urlcolor=linkblue,
    filecolor=magenta,
    pdfauthor={Juan Sebastian Rojas Rodriguez},
    pdftitle={Ultrafast laser parameters optimization through machine learning techniques},
    pdfsubject={Master Thesis},
    pdfkeywords={Quantum Optics, Ultrafast optics, Machine Learning}
}

% --- MINI-ÍNDICE POR CAPÍTULO (estilo Del Valle) ---
% Se carga DESPUÉS de hyperref para que el mini-índice tenga hipervínculos.
% Requiere \dominitoc antes de \tableofcontents (ver main.tex) y
% \chapterminitoc justo tras cada \chapter.
\usepackage{minitoc}
\setcounter{minitocdepth}{1}              % solo secciones (1.1, 1.2, ...), como Del Valle
\mtcsettitle{minitoc}{Contenido}          % encabezado del mini-índice
\mtcsettitlefont{minitoc}{\large\bfseries}% encabezado en negrita
\setlength{\mtcindent}{0pt}               % sin sangría extra a la izquierda
\renewcommand{\mtcSSfont}{\small\rmfamily}% tamaño de las entradas
% \chapterminitoc = mini-índice + un poco de aire antes del texto del capítulo.
\newcommand{\chapterminitoc}{\minitoc\vspace{1.2em}}

% --- INTERLINEADO ---
% \onehalfspacing

% Permite que las páginas no terminen exactamente alineadas abajo
\raggedbottom